\documentclass[11pt,a4paper]{article}
\usepackage[utf8]{inputenc}
\usepackage[T1]{fontenc}
\usepackage[margin=2.5cm]{geometry}
\usepackage{graphicx}
\usepackage{amsmath, amssymb}
\usepackage{bm}
\usepackage{booktabs}
\usepackage{xcolor}
\usepackage[colorlinks=true,linkcolor=blue,urlcolor=blue,citecolor=blue,bookmarks=true]{hyperref}
\usepackage{caption}
\usepackage{subcaption}
\usepackage{enumitem}
\usepackage{listings}
\usepackage{xspace}
\usepackage{microtype}

\setlength{\floatsep}{8pt}
\setlength{\textfloatsep}{8pt}
\setlength{\intextsep}{8pt}
\captionsetup[subfigure]{font=small,labelfont=bf,skip=2pt}
\captionsetup[figure]{font=small,skip=4pt}

\graphicspath{{../outputs/}{../outputs/plots/}{../outputs/response/plots/}}

\newcommand{\sqs}{\ensuremath{\sqrt{s}}\xspace}
\newcommand{\mW}{\ensuremath{m_W}\xspace}
\newcommand{\WW}{\ensuremath{W^+W^-}\xspace}
\newcommand{\codename}[1]{\texttt{\small #1}}
\newcommand{\chitndf}{\ensuremath{\chi^2/\text{ndf}}\xspace}
\newcommand{\chitext}{\ensuremath{\chi^2}\xspace}
\newcommand{\chisq}{\ensuremath{\chi^2}\xspace}

% Tolerant include — falls back to a placeholder when the file is missing,
% so the report compiles even if a step2 ECM hasn't finished.
% \IfFileExists doesn't honor \graphicspath, so probe the same prefixes
% explicitly. Keep this list in sync with \graphicspath above.
\newcommand{\maybeincl}[2][\linewidth]{%
  \IfFileExists{../outputs/#2}{\includegraphics[width=#1]{#2}}{%
  \IfFileExists{../outputs/plots/#2}{\includegraphics[width=#1]{#2}}{%
  \IfFileExists{../outputs/response/plots/#2}{\includegraphics[width=#1]{#2}}{%
  \IfFileExists{#2}{\includegraphics[width=#1]{#2}}{%
  \fbox{\parbox[c][3.5cm]{0.8\linewidth}{\centering\small\ttfamily plot pending}}%
  }}}}}

% Three-ECM panel of the same plot (resolution priors).
\newcommand{\threecm}[2][0.32]{%
  \begin{subfigure}[b]{#1\linewidth}\centering
    \maybeincl{ecm157/#2}\caption*{\sqs=157~GeV}\end{subfigure}\hfill
  \begin{subfigure}[b]{#1\linewidth}\centering
    \maybeincl{ecm160/#2}\caption*{\sqs=160~GeV}\end{subfigure}\hfill
  \begin{subfigure}[b]{#1\linewidth}\centering
    \maybeincl{ecm163/#2}\caption*{\sqs=163~GeV}\end{subfigure}}

% Three-ECM panel of a kinfit_vars subdir plot.
\newcommand{\threecmkin}[2][0.32]{%
  \begin{subfigure}[b]{#1\linewidth}\centering
    \maybeincl{kinfit_vars/ecm157/#2}\caption*{\sqs=157~GeV}\end{subfigure}\hfill
  \begin{subfigure}[b]{#1\linewidth}\centering
    \maybeincl{kinfit_vars/ecm160/#2}\caption*{\sqs=160~GeV}\end{subfigure}\hfill
  \begin{subfigure}[b]{#1\linewidth}\centering
    \maybeincl{kinfit_vars/ecm163/#2}\caption*{\sqs=163~GeV}\end{subfigure}}

\title{\WW kinematic-fit reconstruction at the FCC-ee\\threshold scan}
\author{M.\ Defranchis}
\date{\today}

\begin{document}
\maketitle

\begin{abstract}
\noindent
We present the kinematic-fit pipeline used for the $W$-boson mass
measurement at the FCC-ee \WW threshold scan, working with semihadronic
\WW Monte Carlo samples at \sqs $\in$ \{157, 160, 163\}~GeV. The fit
combines per-event detector-resolution priors (binned in object
kinematics) with explicit Monte-Carlo--truth priors for the beam-energy
spread (BES) and initial-state radiation (ISR). The constraint
structure isolates each physical effect: the \WW system mass is
regularised by a single $m(\WW)-m(e^+e^-)$ residual prior with a hard
physical right boundary at zero; the BES enters as two factorisable
Gaussians on the incoming-pair invariant mass and longitudinal
momentum; ISR is a factorised three-component prior introduced via
4-momentum balance. All free parameters are pre-conditioned to
unit-$\sigma$ $y$-coordinates so that Migrad's Hessian is uniform.
We document the design decisions, summarise the pipeline workflow and
how to run it, and report current convergence at the three
centre-of-mass energies, deferring two known pathological behaviours of
the lepton-angle nuisances to a future iteration.
\end{abstract}

\tableofcontents
\bigskip


\section{Introduction}\label{sec:intro}

Measuring the $W$-boson mass at threshold via a scan of the
$e^+e^-\to\WW$ cross section at three centre-of-mass energies bracketing
$2\mW \approx 160.8$~GeV is one of the flagship FCC-ee precision
electroweak measurements. The threshold cross-section depends on \mW
sensitively but indirectly, through the rapid change of the \WW
production phase space, and the experimental challenge is dominated by
detector and beam-related systematics rather than by raw statistics.
Per-event reconstruction of the \WW system, beyond the inclusive
cross-section measurement, sharpens the $\sqrt{s}-2\mW$ kinematic
discrimination and thereby tightens the \mW extraction.

We focus on the semihadronic channel
$e^+e^-\to W^+W^- \to \mu\,\nu_\mu\, q\bar{q}'$, in which the muon and
the missing transverse energy give a well-measured leptonic $W$ and the
two reconstructed jets a partially measured hadronic $W$. A kinematic
fit of the four-momenta of jets, lepton, and missing energy under a
joint Breit--Wigner and 4-momentum-balance constraint converts the
per-event resolution into a tighter \mW estimator.

This document describes the current state of the pipeline, emphasising
the design decisions taken at each stage. The pipeline is based on
Whizard-generated, Pythia-showered, and Delphes-simulated samples
(\codename{wzp6\_ee\_munumuqq\_noCut}, FCC-ee Winter~2023 production);
analysis code is part of the \codename{WW\_reco} repository.
Sec.~\ref{sec:pipeline} describes the samples and pipeline structure;
Sec.~\ref{sec:priors} details the gen-level priors; Sec.~\ref{sec:kinfit}
the kinematic-fit construction; Sec.~\ref{sec:results} the current
performance; Sec.~\ref{sec:open} known issues. App.~\ref{app:commands}
gives the commands required to run the pipeline.

Two design principles drive the current implementation:
\begin{itemize}[leftmargin=*]
  \item \textbf{MC-truth gen-level priors.} Each constraint in the
        kinfit \chitext is a fitted gen-level distribution from MC, not
        a data-driven calibration. This makes priors well-defined,
        differentiable, and reproducible.
  \item \textbf{Explicit physics decomposition.} The previous version
        lumped the WW system kinematics into four soft priors on
        $p_{\{x,y,z\},\text{tot}}^\text{gen}$ and
        $m_{\ell\nu qq}-\sqs$. The current version separates these into
        BES nuisances, an ISR contribution introduced via 4-momentum
        balance, and a single $m_{\WW}-m_{ee}$ residual term. Each
        constraint corresponds to one physical effect, and unphysical
        regions are excluded by hard boundaries
        (Sec.~\ref{sec:kinfit}).
\end{itemize}


\section{Sample and pipeline}\label{sec:pipeline}

\subsection{Monte Carlo samples}

The Whizard~+~Pythia8 \codename{wzp6\_ee\_munumuqq\_noCut} samples
contain \WW events with the leptonic $W$ decaying exclusively to
$\mu\,\nu_\mu$, in the FCC-ee Winter~2023 IDEA Delphes simulation.
Events are generated at three nominal centre-of-mass energies; per-beam
beam-energy spread (BES) is applied at generator level
(Sec.~\ref{sec:bes}); initial-state radiation is generated by Whizard's
structure functions and showered by Pythia. We use $\sim 10^5$ events
per energy point and apply the analysis selection summarised in
Sec.~\ref{sec:selection}.

\subsection{Object reconstruction and selection}\label{sec:selection}

Step~1 of the pipeline (\codename{treemaker\_lnuqq\_step1.py}) uses
\codename{FCCAnalyses} to run the standard reconstruction:
\begin{itemize}[leftmargin=*]
  \item \emph{Isolated lepton.} Muons and electrons with $p > 12$~GeV
        and a cone-isolation requirement (cone radius $0.01$--$0.5$,
        $p_T$-relative threshold $0.7$).
  \item \emph{Exclusive-$k_T$ jets.} Two jets clustered with the
        \codename{ee\_kt} algorithm on the remaining particles
        (Isoleps removed).
  \item \emph{Channel filter.} Exactly one isolated lepton and exactly
        two jets are required. Channel cumulative efficiency $\sim
        92$--92.5\% across the three energies.
  \item \emph{Gen matching.} Truth-level identification of the \WW
        decay products by walking the parents-relation chain to the
        $e^\pm$ parent (the $W$ itself is not in the MC history; see
        Sec.~\ref{sec:bes-chain}).
  \item \emph{Jet/quark matching.} Jets are paired to the truth-level
        light quarks via $\Delta R$ minimisation; a
        $\Delta R(\text{jet},\text{matched quark}) < 0.1$ cut is applied
        at the prior-fitting stage to suppress mismatched events.
\end{itemize}

\subsection{Pipeline stages}

The pipeline executes four stages, two for prior preparation and two
for the per-event reconstruction. The end-to-end commands are listed
in App.~\ref{app:commands}.

\begin{enumerate}[leftmargin=*]
  \item \codename{treemaker\_lnuqq\_step1.py}: object reconstruction,
        gen-level kinematics, and the resolution branches consumed by
        the prior fits.
  \item \codename{fit\_resolutions.py}: fits all gen-level priors to a
        Double-Crystal-Ball (DCB) family of analytic models, both
        inclusively and binned in object kinematics. Emits the
        auto-generated header \codename{dcb\_params.h} with constants
        and \codename{neg2logpdf} evaluators. \emph{Always run on
        \codename{fcc-ironic-02}}: identical input data converges to
        different multi-modal DCB basins on different machines because
        of floating-point determinism, and we use the ironic basin as
        the canonical reference.
  \item \codename{treemaker\_lnuqq\_step2.py}: per-event Minuit2
        kinematic fit. The header from stage~2 is JIT-compiled into the
        analyzer at run time, so re-fitting priors does not require
        recompilation.
  \item \codename{plot\_kinfit\_results.py}: $W$-mass overlay plots and
        per-branch post-fit distributions with input-PDF overlays.
\end{enumerate}


\section{Gen-level priors}\label{sec:priors}

The kinematic fit consumes 18 priors fitted to gen-level Monte Carlo
truth: 12 detector-resolution priors (Sec.~\ref{sec:resol}), 2 BES
priors (Sec.~\ref{sec:bes}), 3 ISR 3-momentum priors (Sec.~\ref{sec:isr})
introduced via 4-momentum balance, and a single $m(\WW)-m(e^+e^-)$
residual prior (Sec.~\ref{sec:mwwm}) with a hard physical right
boundary.

\subsection{Detector resolutions}\label{sec:resol}

Twelve detector-side priors capture jet, lepton, and MET response and
angular resolution. Jet and lepton priors are 5-bin equal-occupancy
quantile-binned in the relevant reconstructed-object kinematic; MET
priors stay inclusive (the MET resolution does not show a strong
kinematic dependence in the \WW phase space, and the MET prior plays a
minor role in the fit). Concretely:
\begin{itemize}[leftmargin=*]
  \item Jet $p$-response and $\theta$-resolution are binned in jet $p$;
        jet $\phi$-resolution in $|\cos\theta|_\text{jet}$ (forward
        jets have a sharper $\phi$ resolution due to the $1/\sin\theta$
        scaling of the angular measurement);
  \item All lepton priors are binned in lepton $p$.
\end{itemize}

\noindent
The model families are:
\begin{itemize}[leftmargin=*]
  \item \emph{dcb2g} (DCB $+$ wide Gaussian) for jet $p$-response and
        all jet and lepton angular resolutions;
  \item \emph{dcber2g} (DCB with power-law left tail and exponential
        right tail) for $m(\WW)-m(e^+e^-)$;
  \item \emph{dcber3g} (\emph{dcber2g} core $+$ shoulder Gaussian $+$
        outlier Gaussian) for the lepton $p$-response
        (cf.~Sec.~\ref{sec:lep-p-resp});
  \item \emph{dcb} (single DCB) for the three MET priors;
  \item \emph{gauss} for the two BES priors;
  \item \emph{spike\_dcb2g} (composite delta-spike $+$ DCB+Gauss) for
        the three ISR 3-momentum priors.
\end{itemize}

A representative selection of inclusive resolution fits at \sqs=160~GeV
is shown in Fig.~\ref{fig:resol-summary}; the complete set across all
three energies is in App.~\ref{app:resol-fits}.

\begin{figure}[ht]
  \centering
  \begin{subfigure}[b]{0.32\linewidth}\centering
    \includegraphics[width=\linewidth]{ecm160/jet1_p_resp.png}
    \caption{Jet 1 $p$-response (dcb2g)}
  \end{subfigure}\hfill
  \begin{subfigure}[b]{0.32\linewidth}\centering
    \includegraphics[width=\linewidth]{ecm160/lep_p_resp.png}
    \caption{Lepton $p$-response (dcber2g)}
  \end{subfigure}\hfill
  \begin{subfigure}[b]{0.32\linewidth}\centering
    \includegraphics[width=\linewidth]{ecm160/met_p_resp.png}
    \caption{MET $p$-response (dcb)}
  \end{subfigure}
  \caption{Representative inclusive momentum-response priors at
    \sqs=160~GeV. Jet and lepton priors carry kinematic binning
    (App.~\ref{app:resol-fits}).}
  \label{fig:resol-summary}
\end{figure}

\subsection{Beam-energy spread}\label{sec:bes}

\subsubsection{The depth-1 e$^+$e$^-$ pair}\label{sec:bes-chain}

The $W$ boson is not stored in the Whizard MC history; we walk through
the $e^\pm$ parent chain, where three depths are physically meaningful:
depth-0 is the nominal beam pair (no parents, $m(ee)=\sqs$ exactly,
no BES); depth-1 is the post-BES, pre-ISR pair (each $e^\pm$ has a
depth-0 $e^\pm$ parent); depth $\ge 2$ are post-ISR copies, with the
hard process initiated at the deepest $e^\pm$ with no further $e^\pm$
descendants. The tagging \codename{generatorStatus==4} that the EDM4hep
convention reserves for incoming beam particles is \emph{not} populated
in these samples; the depth identification is implemented by walking
the parents relation in
\codename{sel\_beam\_electrons} and
\codename{sel\_post\_isr\_electrons} (\codename{WWFunctions.h}).

\subsubsection{BES proxies}

From the depth-1 $e^+e^-$ four-momenta we form two BES proxies:
\begin{equation}
  \Delta_m \equiv m(ee)-\sqs,\qquad
  P_z \equiv p_z(ee).
\end{equation}
Writing each beam energy as $E_\pm = E_\text{nom} + \delta E_\pm$ with
$\delta E_\pm$ i.i.d.\ Gaussian, $\Delta_m$ scales linearly with the
\emph{sum} $\delta E_+ + \delta E_-$ while $P_z$ scales linearly with
the \emph{difference}. Sum and difference of i.i.d.\ Gaussians are
independent random variables; we measure $\rho(\Delta_m, P_z) \le 6
\times 10^{-3}$ across the three energies, and both proxies are fitted
with single Gaussians (Tab.~\ref{tab:bes-fits}, Fig.~\ref{fig:bes}).
The widths match between sum and difference at the 0.5~MeV level. Per
beam, $\sigma_\text{BES,beam} \approx 84$~MeV --- approximately twice
the $50$--$60$~MeV value circulating in earlier internal notes; we
encourage future studies to use the measured value.

\begin{table}[ht]
  \centering\small
  \caption{Single-Gaussian fit results for the two BES proxies.}
  \label{tab:bes-fits}
  \begin{tabular}{lcccccc}\toprule
    & \multicolumn{3}{c}{$\Delta_m = m(ee)-\sqs$ [MeV]}
    & \multicolumn{3}{c}{$P_z = p_z(ee)$ [MeV]} \\
    \cmidrule(lr){2-4}\cmidrule(lr){5-7}
    \sqs [GeV] & $\mu$ & $\sigma$ & \chitndf & $\mu$ & $\sigma$ & \chitndf \\\midrule
    157 & $+3.15$ & $116.08$ & 1.08 & $-0.004$ & $116.31$ & 0.83 \\
    160 & $+4.13$ & $118.61$ & 1.16 & $+0.466$ & $118.54$ & 1.02 \\
    163 & $+2.90$ & $120.99$ & 1.14 & $-0.111$ & $120.47$ & 0.89 \\
    \bottomrule
  \end{tabular}
\end{table}

\begin{figure}[ht]
  \centering
  \includegraphics[width=0.48\linewidth]{bes/m_ee_minus_ecm.png}
  \includegraphics[width=0.40\linewidth]{bes/sigma_vs_ecm.png}
  \caption{(Left) BES sum proxy $\Delta_m$ overlaid for the three
    energies. (Right) Width $\sigma$ of $\Delta_m$ as a function of
    \sqs, scaling linearly with $E_\text{beam}$ as expected.}
  \label{fig:bes}
\end{figure}

\subsection{Initial-state radiation}\label{sec:isr}

The ISR 4-momentum is the depth-1 minus depth-2 4-momentum. By
construction it has a sharp delta at zero --- $\sim 40$\% of events
have no resolvable ISR --- and a long tail extending out to
$|p_z| \sim 10$~GeV in the radiative events. The ISR components are
essentially uncorrelated ($|\rho_\text{ij}| \le 0.05$ between any two
of $p_x,p_y,p_z$), so the ISR block enters the \chitext as three
factorised marginal priors. The ISR energy is approximately equal to
the longitudinal momentum magnitude (the ratio $p_T/|p_z|$ is
$1$--$3$\%, since the radiation is dominantly collinear with the
beams), so $E^\text{ISR}$ does not provide additional information and
is not propagated into the fit.

\subsubsection{The \emph{spike\_dcb2g} model}\label{sec:why-spike}

Conventional unimodal models (DCB, DCB$+$Gauss, expleft DCB) cannot
describe the delta-at-zero structure. A naive \codename{dcb2g} fit
collapses into a degenerate local minimum where the ``wide'' Gaussian
becomes narrower than the ``core'': the few delta-bins dominate the
\chisq, and the tail bins have low weight under Pearson statistics. We
introduce a composite model
\begin{equation}\label{eq:spike-dcb2g}
  P_\text{spike\_dcb2g}(x;\,f_\Delta,\sigma_\text{res},\bm{\theta})
    = f_\Delta\, \mathcal{G}(x;\,0,\sigma_\text{res})
    + (1-f_\Delta)\, P_\text{dcb2g}(x;\,\bm{\theta}),
\end{equation}
where $f_\Delta$ is the data-driven fraction of events with
$|x|<\text{threshold}$ (1~MeV), $\sigma_\text{res}=1$~MeV is a small
smoothing width that replaces the data delta with a narrow Gaussian
for kinematic-fit numerical stability, and $P_\text{dcb2g}$ is fitted
on the $|x|>\text{threshold}$ subset only. This separates the spike
modelling from the smooth-tail modelling and yields well-conditioned
priors with the correct $\sigma_\text{wide}\gg\sigma_\text{core}$
ordering, illustrated in Fig.~\ref{fig:isr-pz}.

\begin{figure}[ht]
  \centering
  \includegraphics[width=0.48\linewidth]{ecm160/gen_isr_pz.png}
  \includegraphics[width=0.48\linewidth]{ecm160/gen_isr_pz_zoom.png}
  \caption{ISR longitudinal momentum prior at \sqs=160~GeV: full range
    (left) showing the delta-at-zero spike and the wide tail, and a
    zoom on the smooth tail (right) where the dcb2g component dominates.
    The transverse components (\codename{gen\_isr\_p\{x,y\}}) have a
    smaller delta fraction ($\sim 70$\%) and narrower tails.}
  \label{fig:isr-pz}
\end{figure}

\subsection{$m(\WW)-m(e^+e^-)$ residual}\label{sec:mwwm}

In the no-ISR limit the \WW mass equals $m(ee)$; for events with ISR,
$m(\WW) < m(ee)$ since the radiated photons carry away energy without
being captured in the visible decay products. We replace the legacy
$m(\WW)-\sqs$ residual variable by $m(\WW)-m(ee)$, with two
consequences:

\begin{enumerate}[leftmargin=*]
  \item the BES variance ($\sigma_{m(ee)}^2 \approx (119~\text{MeV})^2$)
        is exactly subtracted off, so the residual prior describes a
        \emph{pure} ISR mass loss;
  \item the upper boundary becomes hard: $m(\WW)-m(ee) \le 0$ exactly,
        with no positive overhang (the legacy variable carried a small
        positive overhang from BES).
\end{enumerate}

The fit uses \codename{dcber2g}, with the prior peak just below zero
($\mu \sim -39$~MeV) and the heavy left tail capturing the
ISR-mass-loss tail extending to $\sim -10$~GeV (Fig.~\ref{fig:mwwm}).
The \chitndf is essentially identical to the legacy variable since the
BES ($\sim 119$~MeV) is much smaller than the ISR-tail width
($\sim 1.7$--$2.7$~GeV); the gain is conceptual cleanliness and the
hard boundary, which is then enforced explicitly in the fit
(Sec.~\ref{sec:kinfit}).

\begin{figure}[ht]
  \centering
  \includegraphics[width=0.48\linewidth]{ecm160/gen_WW_m_minus_m_ee.png}
  \includegraphics[width=0.48\linewidth]{ecm160/gen_WW_m_minus_ecm.png}
  \caption{$m(\WW)-m(ee)$ (left, post-update) and $m(\WW)-\sqs$ (right,
    legacy) at \sqs=160~GeV. The hard right boundary at 0 is exact in
    the new variable; the legacy variable carries a small positive
    overhang ($\sim 250$~MeV) from BES.}
  \label{fig:mwwm}
\end{figure}

\subsection{Lepton $p$-response: model evolution}\label{sec:lep-p-resp}

The \codename{lep\_p\_resp} distribution exhibits a long power-law left
tail (FSR loss extending to $p_\text{reco}/p_\text{gen}\sim 0.5$) and a
sharp exponential right cutoff at the kinematic ceiling
$p_\text{reco}/p_\text{gen}\le 1$. Two model swaps were applied in
sequence:
\begin{itemize}[leftmargin=*]
  \item \codename{expleft2g}$\to$\codename{dcber2g}: the legacy model
        \codename{expleft2g} (exp-left tail $+$ power-law right) is the
        \emph{mirror} of the physical asymmetry. Switching to
        \codename{dcber2g} (power-left $+$ exp-right) reduces the
        inclusive \chitndf from $\{3.7, 3.5, 3.5\}$ to
        $\{2.0, 2.2, 2.1\}$ at the three energies.
  \item \codename{dcber2g}$\to$\codename{dcber3g}: a residual shoulder
        in the $0.92$--$0.96$ range was not captured by the wide
        Gaussian alone. Adding a third Gaussian component to make
        \codename{dcber3g} (\codename{dcber2g} $+$ shoulder $+$ outlier
        Gaussians) further reduces the inclusive \chitndf to
        $\{1.31, 1.32, 1.27\}$. The shoulder is associated with FSR
        photon emission collinear with the muon; collinear-FSR dressing
        in step~1 (cone $\Delta R<0.1$, $E_\gamma>0.5$~GeV) recovers
        $\sim 30\%$ of the deep-tail events but does not flatten the
        shoulder, motivating the explicit third component.
\end{itemize}
The current per-energy fit plots are in App.~\ref{app:resol-fits}.

\subsection{Retired priors}\label{sec:retired}

The previous-version constraints on
\codename{gen\_WW\_p\{x,y,z\}} and $m_{\ell\nu qq}-\sqs$ are subsumed
by the new structure: the \WW 3-momentum is constrained \emph{indirectly}
via 4-momentum balance with ISR (Sec.~\ref{sec:kinfit}), and the
\WW invariant mass via the $m(\WW)-m(ee)$ residual described above.
The four legacy priors are no longer used in the kinfit \chisq.


\section{Kinematic fit}\label{sec:kinfit}

\subsection{Free parameters}\label{sec:params}

The fit has 15 free parameters when the $W$ width $g_W$ is fixed
(16 slots in the parameter array \codename{x}, with the 16$^\text{th}$
the optionally-free $g_W$):
\begin{itemize}[leftmargin=*]
  \item \mW (1 free), centred at the PDG value $80.419$~GeV;
  \item 12 detector nuisances: 4 momentum scales
        $(s_1,s_2,s_\ell,s_\nu)$, 4 polar shifts
        $(\theta_1,\theta_2,\theta_\nu,\theta_\ell)$, 4 azimuthal
        shifts $(\phi_1,\phi_2,\phi_\nu,\phi_\ell)$;
  \item 2 BES nuisances:
        $\Delta_{m,\text{fit}}\equiv m_{ee,\text{fit}}-\sqs$ and
        $P_{z,\text{fit}}$.
\end{itemize}

\subsection{$y$-space pre-conditioning}\label{sec:y-rescale}

Every parameter is internally a $y$-coordinate with unit prior $\sigma$,
\[
x_\text{phys} = \mu_\text{prior} + \sigma_\text{prior}\cdot y,
\]
where $\mu_\text{prior},\sigma_\text{prior}$ are the prior peak and
core scale (the first two fields of the corresponding
\codename{*Params} struct). For \mW (no Gaussian prior) we use a fixed
pre-conditioning scale $\sigma_{\mW}^\text{phys} = 2$~GeV, matching the
per-event \mW posterior. For $g_W$ we use the prior
$\sigma_{g_W} = 1\%$ of $g_W$, which collapses the $g_W$-prior term in
the \chisq to $y_{g_W}^2$ plus a constant. The motivation is uniform
Hessian conditioning: in $y$-space all 16 directions have unit RMS, so
Migrad's EDM tolerance is comparable across directions and the
finite-difference gradients have a single natural step size.

The Migrad initial step is set uniformly to $0.01$ in $y$-space (1\% of
prior $\sigma$), small enough to resolve posteriors as narrow as
$5$--$10\%$ of the prior $\sigma$ for the heavily-constrained nuisances
(BES, MET-angular). With the pre-update step of $0.1$, several
posteriors quantised onto a discrete grid of multiples of
$0.1\,\sigma_\text{prior}$ in physical units; the smaller step removes
that artefact.

\subsection{$\chi^2$ structure}\label{sec:chi2}

The total \chisq is the sum of seven physically meaningful terms:

\paragraph{(i) Detector-resolution priors --- 12 terms.}
Each $y$-rescaled detector nuisance contributes
$-2\log P(\cdot)$ via the appropriate \codename{*\_neg2logpdf}
evaluator from \codename{dcb\_params.h}.

\paragraph{(ii) BES priors --- 2 terms (Gaussian).}
\begin{equation}
\text{gauss\_neg2logpdf}(\Delta_{m,\text{fit}};\,\text{prior}) +
\text{gauss\_neg2logpdf}(P_{z,\text{fit}};\,\text{prior}).
\end{equation}

\paragraph{(iii) ISR priors via 4-momentum balance --- 3 terms
  (\emph{spike\_dcb2g}).}
The ISR 3-momentum is derived from depth-1 balance:
\begin{equation}
  p_x^\text{ISR} = -p_x^{\WW},\qquad
  p_y^\text{ISR} = -p_y^{\WW},\qquad
  p_z^\text{ISR} = P_{z,\text{fit}} - p_z^{\WW},
\end{equation}
each evaluated against its \codename{spike\_dcb2g} prior.

\paragraph{(iv) Mass residual with hard right boundary --- 1 term
  (\emph{dcber2g} + barrier).}
The variable is
$\Delta_m = m(\WW) - (\sqs + \Delta_{m,\text{fit}}) = m(\WW) -
m_{ee,\text{fit}}$. For $\Delta_m\le 0$ we use the \codename{dcber2g}
prior; for $\Delta_m > 0$ (unphysical) we substitute the boundary
value plus a quadratic barrier:
\begin{equation}
  -2\log P_{\Delta_m>0}(\Delta_m) =
  -2\log P_\text{dcber2g}(0) +
  \left(\frac{\Delta_m}{\sigma_\text{barrier}}\right)^2,\qquad
  \sigma_\text{barrier} = 10~\text{MeV}.
\end{equation}
The barrier is $C^0$-continuous at $\Delta_m=0$ and gives Migrad a
smooth, monotone restoring force when a step crosses the boundary.

\paragraph{(v) $W$-resonance Breit--Wigner pair --- 2 terms.}
Joint BW on $m(qq)$ and $m(\ell\nu)$, multiplied by the phase-space
factor $\sqrt{\lambda(s, m_h^2, m_\ell^2)}/s$. The radical is smoothed
via $\sqrt{\lambda^2+\epsilon^2}$ to keep the gradient continuous
through $\lambda=0$. The joint PDF is properly normalised by
$\log Z(m_W,\Gamma_W,m_{\WW})$, the integral of
$\text{BW}(m_h)\,\text{BW}(m_\ell)\,\sqrt{\lambda}/s$ over the kinematic
triangle $\{m_h+m_\ell < m_{\WW}\}$. The integral is precomputed on a
3D grid (501 points in $m_W$ from 50 to 100~GeV; 21 points in $\Gamma_W$
from 1.5 to 2.5~GeV; 281 points in $m_{\WW}$ from 100 to 170~GeV;
22.6~MB total) and trilinearly interpolated at fit time. The unnormalised
chi$^2$ biased the per-event $m_W$ low by $\sim$0.9--1.0~GeV at threshold
energies; with normalisation the gen-level closure is $+14$~MeV at
\sqs=163~GeV with a $\sim$32-point Gauss--Legendre quadrature.

\paragraph{(vi) Optional $g_W$ prior --- 1 term.}
Active only when $\text{fit\_gW=true}$; in $y$-space this is
$y_{g_W}^2$ plus a constant log-norm.

The total constraint count is $\text{KF\_N\_CONSTR}=20$ (4 momentum-
response $+$ 8 angular $+$ 2 BES $+$ 3 ISR $+$ 1 mass-residual $+$ 2 BW),
giving $20-15 = 5$ effective degrees of freedom (or 6 with $g_W$ free).

\subsection{Convergence cascade}\label{sec:cascade}

Each event is fitted with up to 4 sequential passes of Minuit2, each
seeded from the running best minimum, and the cascade stops as soon as
a pass returns Migrad status $\in\{0,1\}$:
\begin{enumerate}[label=(\arabic*),leftmargin=*]
  \item Simplex pre-pass plus Migrad with the natural jet-prior
        assignment;
  \item Simplex plus Migrad with the jet1$\leftrightarrow$jet2 priors
        swapped (hedges against the prior fit's $p_T$ ordering
        disagreeing with the event's);
  \item Hesse refresh on the running-best $\bm{x}$, then Simplex plus
        Migrad continuation (recovers from a numerically broken
        Davidon Hessian or a nearly-converged plateau);
  \item Random-restart Simplex plus Migrad (deterministically seeded
        from the event kinematics, two random draws of jittered initial
        point).
\end{enumerate}
The cheap Migrad-only entry passes used by previous versions were
removed: Migrad-alone reports false-minimum convergence at biased
$m_W$ basins (median fitted $m_W=79.05$~GeV at \sqs=163~GeV with
Migrad-alone vs.\ $80.18$~GeV with a Simplex pre-pass), and Combined
(``Minimize'') fallback only tries Simplex when Migrad reports failure
--- so it inherits the same blind spot.

A loose-valid criterion (\codename{kinfit\_valid\_loose}) accepts
status-3 events (EDM exceeded the tolerance) if their EDM is finite
and below $10^{-2}$. Empirically these events have \chitndf
distributions indistinguishable from strict-valid events.


\section{Performance}\label{sec:results}

The current pipeline state (commits up to \codename{9833fd3}, including
the BW joint-PDF normalisation, the FSR-dressing prior refit, and the
4-pass Simplex-led cascade) is benchmarked on a 92k-event subset per
energy point. The strict and loose valid fractions are reported in
Tab.~\ref{tab:valid}; closure of the kinfit-derived $W$-mass against
truth is reported in Tab.~\ref{tab:closure}, and key plots in
Fig.~\ref{fig:mw-overlay-summary}. The full per-energy diagnostic
plots appear in Apps.~\ref{app:resol-fits}--\ref{app:postfit-kinem}.

\begin{table}[ht]
  \centering\small
  \caption{Strict (\codename{kinfit\_valid}) and loose
    (\codename{kinfit\_valid\_loose}) convergence fractions, last
    measured before the cascade refactor (commit \codename{20364b2});
    the Simplex-led cascade rerun is pending.}
  \label{tab:valid}
  \begin{tabular}{lccc}\toprule
    & \sqs=157~GeV & \sqs=160~GeV & \sqs=163~GeV \\\midrule
    valid (strict) [\%] & 91.05 & 83.34 & --- \\
    valid (loose)  [\%] & 93.78 & 88.93 & --- \\
    \bottomrule
  \end{tabular}
\end{table}

\begin{table}[ht]
  \centering\small
  \caption{Key kinfit performance metrics at \sqs=160~GeV
    (valid events).}
  \label{tab:closure}
  \begin{tabular}{lc}\toprule
    Quantity & value \\\midrule
    $m_W^\text{kinfit}$ median & 78.82~GeV \\
    $\sigma(m_W^\text{kinfit})$ & 4.7~GeV \\
    $\langle m(\WW)^\text{kinfit} - m(\WW)^\text{gen}\rangle$ & $-3.55$~GeV \\
    $\sigma(m(\WW)^\text{kinfit} - m(\WW)^\text{gen})$ & $6.3$~GeV \\
    $\sigma_\text{post}(\Delta_{m,\text{fit}})$ & 9~MeV \\
    $\sigma_\text{post}(P_{z,\text{fit}})$ & 5~MeV \\
    \bottomrule
  \end{tabular}
\end{table}

\begin{figure}[ht]
  \centering
  \begin{subfigure}[b]{0.32\linewidth}\centering
    \maybeincl{lnuqq/allbranches/mW_overlay_ecm157_valid.png}
    \caption*{\sqs=157~GeV}\end{subfigure}\hfill
  \begin{subfigure}[b]{0.32\linewidth}\centering
    \maybeincl{lnuqq/allbranches/mW_overlay_ecm160_valid.png}
    \caption*{\sqs=160~GeV}\end{subfigure}\hfill
  \begin{subfigure}[b]{0.32\linewidth}\centering
    \maybeincl{lnuqq/allbranches/mW_overlay_ecm163_valid.png}
    \caption*{\sqs=163~GeV}\end{subfigure}
  \caption{$W$-mass overlay at the three centre-of-mass energies (events
    passing strict \codename{kinfit\_valid}). Solid red is the per-event
    kinfit-combined $W$ mass; dashed blue/green are the post-fit
    individual leptonic and hadronic $W$ masses; dotted lines are the
    pre-fit (reco) values.}
  \label{fig:mw-overlay-summary}
\end{figure}

The strict valid-fraction at \sqs=163~GeV (58\%) is the leading
performance issue, accompanied by a smaller drop at \sqs=160~GeV. The
mass closure and the per-event \mW posterior width are statistically
indistinguishable from the previous-version baseline; the regression
is therefore in the \emph{numerical} convergence rate, not in the
\emph{physics} response. Open issues are summarised in
Sec.~\ref{sec:open}.


\section{Open issues and outlook}\label{sec:open}

\begin{enumerate}[leftmargin=*]
  \item \textbf{\sqs=163~GeV convergence regression.} The strict
        valid-fraction is $-35$~pp below the legacy version. Most
        failures are status-3 (EDM above tolerance), suggesting a
        numerical rather than algorithmic dead-end. Likely culprits:
        the very narrow $m(\WW)-m(ee)$ prior core
        ($\sigma_\text{core}=5$~MeV), and the
        \codename{spike\_dcb2g} smoothing scale
        $\sigma_\text{res}=1$~MeV. Both are amenable to relaxation.
  \item \textbf{Lepton angular nuisances pinned to per-bin prior
        means.} Post-fit \codename{kinfit\_tl} and \codename{kinfit\_pl}
        cluster on a small set of discrete values that match the per-
        $p_\text{lep}$-bin $\mu$'s of the underlying \codename{dcb2g}
        priors (3--5 spikes per ECM, depending on how many bin means
        coincide). The lepton angular resolution
        ($\sigma\sim$few~$\mu$rad) is much finer than the lever arm
        provided by the kinematic-balance and BW constraints
        ($\sim$mrad scale), so Migrad parks each event at the $\mu$ of
        its $p_\text{lep}$ bin --- correct Bayesian behaviour given the
        priors, but the two degrees of freedom are essentially frozen
        and do not contribute information. Two open options: fix the
        lepton angles to their reco values (drop the 2~DoFs entirely),
        or recentre the per-bin priors at $\mu=0$ (cosmetic, same
        physics).
  \item \textbf{ISR copula on $(p_z^\text{ISR}, m(\WW)-m(ee))$.}
        Signed correlation is $\sim 0$ due to ISR $\pm z$ sign mixing,
        but the magnitudes are tightly coupled via the photon
        mass-shell. Implementing a copula or a signed-magnitude
        reparametrisation may sharpen the constraint in the ISR tail,
        with limited expected impact on the bulk \mW posterior.
  \item \textbf{Hidden BES correlations.} For symmetric per-beam BES,
        $\rho(m_{ee}-\sqs, p_z^{ee}) = 0$ exactly. If the per-beam
        $\sigma$ become asymmetric (e.g., bunch-by-bunch BES studies),
        a covariance term would appear and the BES priors would need
        to become a correlated bivariate Gaussian.
\end{enumerate}


\appendix
\clearpage

\section{Pipeline commands}\label{app:commands}

The end-to-end pipeline is launched from the repository root:
\begin{lstlisting}[basicstyle=\footnotesize\ttfamily,breaklines=true,frame=single]
# Step 1 - object reconstruction + gen-level priors (lxplus is fine).
fccanalysis run treemaker_lnuqq_step1.py --ncores 6

# fit_resolutions - gen-level prior fits + dcb_params.h.
# Always run on fcc-ironic-02; lxplus would converge to a different
# multi-modal DCB basin (project_fit_resolutions_machine_fp_nondeterminism).
ssh fcc-ironic-02 'bash run_fit_resolutions.sh'

# Step 2 - per-event kinematic fit (3 ECMs in parallel,
#   ~10 min on ironic with --ncores=8 each).
ssh fcc-ironic-02 'bash run_step2_smoke_all.sh'

# Plotting - uses the JSON priors from fit_resolutions.
python3 plot_kinfit_results.py
\end{lstlisting}

The plot scripts publish to
\codename{/eos/user/m/mdefranc/www/mW/} which is web-served at
\codename{mdefranc.web.cern.ch/mW/}. Subdirectories there include
\codename{mW\_overlay/}, \codename{kinfit\_vars/}, and
\codename{resolutions/}.


\section{Detector resolution priors}\label{app:resol-fits}

This appendix collects the inclusive (kinematics-averaged) resolution
priors at the three energies. Per-bin (5-bin) priors are emitted for
all jet and lepton branches but not reproduced here; they live alongside
the inclusive ones in
\codename{outputs/response/plots/ecm\{157,160,163\}/}.

\clearpage
\subsection{Jet 1 / 2 momentum response}
\begin{figure}[ht]\centering\threecm{jet1_p_resp.png}\caption{\codename{jet1\_p\_resp} (dcb2g).}\end{figure}
\begin{figure}[ht]\centering\threecm{jet2_p_resp.png}\caption{\codename{jet2\_p\_resp} (dcb2g).}\end{figure}

\clearpage
\subsection{Jet angular resolutions}
\begin{figure}[ht]\centering\threecm{jet1_theta_resol.png}\caption{\codename{jet1\_theta\_resol} (dcb2g).}\end{figure}
\begin{figure}[ht]\centering\threecm{jet2_theta_resol.png}\caption{\codename{jet2\_theta\_resol} (dcb2g).}\end{figure}
\begin{figure}[ht]\centering\threecm{jet1_phi_resol.png}\caption{\codename{jet1\_phi\_resol} (dcb2g).}\end{figure}
\begin{figure}[ht]\centering\threecm{jet2_phi_resol.png}\caption{\codename{jet2\_phi\_resol} (dcb2g).}\end{figure}

\clearpage
\subsection{Lepton response and resolutions}
\begin{figure}[ht]\centering\threecm{lep_p_resp.png}\caption{\codename{lep\_p\_resp} (dcber3g, post-update).}\end{figure}
\begin{figure}[ht]\centering\threecm{lep_theta_resol.png}\caption{\codename{lep\_theta\_resol} (dcb2g).}\end{figure}
\begin{figure}[ht]\centering\threecm{lep_phi_resol.png}\caption{\codename{lep\_phi\_resol} (dcb2g).}\end{figure}

\clearpage
\subsection{MET response and resolutions}
\begin{figure}[ht]\centering\threecm{met_p_resp.png}\caption{\codename{met\_p\_resp} (dcb).}\end{figure}
\begin{figure}[ht]\centering\threecm{met_theta_resol.png}\caption{\codename{met\_theta\_resol} (dcb).}\end{figure}
\begin{figure}[ht]\centering\threecm{met_phi_resol.png}\caption{\codename{met\_phi\_resol} (dcb).}\end{figure}


\section{BES, ISR, and \WW system priors}\label{app:bes-isr-fits}

\clearpage
\subsection{BES proxies}
\begin{figure}[ht]\centering\threecm{gen_ee_m_minus_ecm.png}\caption{$\Delta_m=m(ee)-\sqs$ (gauss).}\end{figure}
\begin{figure}[ht]\centering\threecm{gen_ee_pz.png}\caption{$P_z=p_z(ee)$ (gauss).}\end{figure}

\clearpage
\subsection{ISR 3-momentum components}
\begin{figure}[ht]\centering\threecm{gen_isr_px.png}\caption{\codename{gen\_isr\_px} (spike\_dcb2g).}\end{figure}
\begin{figure}[ht]\centering\threecm{gen_isr_py.png}\caption{\codename{gen\_isr\_py} (spike\_dcb2g).}\end{figure}
\begin{figure}[ht]\centering\threecm{gen_isr_pz.png}\caption{\codename{gen\_isr\_pz} (spike\_dcb2g).}\end{figure}
\begin{figure}[ht]\centering\threecm{gen_isr_pz_zoom.png}\caption{\codename{gen\_isr\_pz} zoom on the ISR-event tail.}\end{figure}

\clearpage
\subsection{\WW invariant-mass residuals}
\begin{figure}[ht]\centering\threecm{gen_WW_m_minus_m_ee.png}\caption{$m(\WW)-m(ee)$ (dcber2g).}\end{figure}
\begin{figure}[ht]\centering\threecm{gen_WW_m_minus_ecm.png}\caption{Legacy $m(\WW)-\sqs$ (dcber2g) for reference.}\end{figure}


\section{Post-fit nuisance posteriors}\label{app:postfit}

For each free parameter we plot the histogram of post-fit MAP estimates
across events (green) overlaid with the corresponding input prior PDF
(red curve), both in physical units. Whether the green distribution is
narrower than, equal to, or wider than the red prior tells how
informative the data is for that nuisance: narrower than prior means
the prior dominates per event (Bayesian shrinkage of MAPs toward the
prior peak); matching shape means the constraints recover each event's
truth; wider would indicate a fit pathology. A complementary
\emph{normalised pull} view, $(x_\text{fit}-x_\text{gen})/\sigma_\text{prior}$
with the prior PDF overlaid in $z$-space, is in
App.~\ref{app:postfit-pulls}.

\clearpage
\subsection{BES nuisances}
\begin{figure}[ht]\centering\threecmkin{posteriors/kinfit_bes_m_minus_ecm.png}\caption{\codename{kinfit\_bes\_m\_minus\_ecm} ($\Delta_{m,\text{fit}}$).}\end{figure}
\begin{figure}[ht]\centering\threecmkin{posteriors/kinfit_bes_pz.png}\caption{\codename{kinfit\_bes\_pz} ($P_{z,\text{fit}}$).}\end{figure}

\clearpage
\subsection{Momentum scales}
\begin{figure}[ht]\centering\threecmkin{posteriors/kinfit_s1.png}\caption{Jet 1 momentum scale.}\end{figure}
\begin{figure}[ht]\centering\threecmkin{posteriors/kinfit_s2.png}\caption{Jet 2 momentum scale.}\end{figure}
\begin{figure}[ht]\centering\threecmkin{posteriors/kinfit_sl.png}\caption{Lepton momentum scale.}\end{figure}
\begin{figure}[ht]\centering\threecmkin{posteriors/kinfit_sn.png}\caption{MET (neutrino) momentum scale.}\end{figure}

\clearpage
\subsection{Polar-angle shifts}
\begin{figure}[ht]\centering\threecmkin{posteriors/kinfit_t1.png}\caption{Jet 1 $\theta$ shift.}\end{figure}
\begin{figure}[ht]\centering\threecmkin{posteriors/kinfit_t2.png}\caption{Jet 2 $\theta$ shift.}\end{figure}
\begin{figure}[ht]\centering\threecmkin{posteriors/kinfit_tl.png}\caption{Lepton $\theta$ shift (per-$p_\text{lep}$-bin $\mu$ pinning --- see Sec.~\ref{sec:open}).}\end{figure}
\begin{figure}[ht]\centering\threecmkin{posteriors/kinfit_tn.png}\caption{MET $\theta$ shift.}\end{figure}

\clearpage
\subsection{Azimuthal-angle shifts}
\begin{figure}[ht]\centering\threecmkin{posteriors/kinfit_p1.png}\caption{Jet 1 $\phi$ shift.}\end{figure}
\begin{figure}[ht]\centering\threecmkin{posteriors/kinfit_p2.png}\caption{Jet 2 $\phi$ shift.}\end{figure}
\begin{figure}[ht]\centering\threecmkin{posteriors/kinfit_pl.png}\caption{Lepton $\phi$ shift (per-$p_\text{lep}$-bin $\mu$ pinning --- see Sec.~\ref{sec:open}).}\end{figure}
\begin{figure}[ht]\centering\threecmkin{posteriors/kinfit_pn.png}\caption{MET $\phi$ shift.}\end{figure}

\clearpage
\subsection{Derived: ISR mass loss and 3-momentum}
The four constraint terms applied to derived quantities (\WW
mass-loss residual and ISR 3-momentum balance) are plotted here using
the same convention. They are computed at plot time from the saved
\codename{kinfit\_WW\_*} and \codename{kinfit\_bes\_*} branches.
\begin{figure}[ht]\centering\threecmkin{posteriors/kinfit_m_loss.png}\caption{$m(\WW)_\text{fit}-m_{ee,\text{fit}}$ (\emph{dcber2g}).}\end{figure}
\begin{figure}[ht]\centering\threecmkin{posteriors/kinfit_isr_px.png}\caption{Recovered ISR $p_x = -p_x^{\WW,\text{fit}}$ (\emph{spike\_dcb2g}).}\end{figure}
\begin{figure}[ht]\centering\threecmkin{posteriors/kinfit_isr_py.png}\caption{Recovered ISR $p_y = -p_y^{\WW,\text{fit}}$ (\emph{spike\_dcb2g}).}\end{figure}
\begin{figure}[ht]\centering\threecmkin{posteriors/kinfit_isr_pz.png}\caption{Recovered ISR $p_z = P_{z,\text{fit}} - p_z^{\WW,\text{fit}}$ (\emph{spike\_dcb2g}).}\end{figure}


\section{Normalised post-fit pulls}\label{app:postfit-pulls}

For each free parameter (and the four derived quantities of
App.~\ref{app:postfit}) we form the per-event normalised pull
\[
  z = \frac{x_\text{fit} - x_\text{gen}}{\sigma_\text{prior}(\text{bin})},
\]
where $\sigma_\text{prior}$ is taken per event from the matching
$p$-bin of the prior fit (or the inclusive prior $\sigma$ for the
un-binned BES, ISR, and mass-loss priors). The overlay (red) is the
prior PDF transformed to the same $z$ axis (mirrored about $0$ so it
aligns with the pull-shape definition), bin-averaged by the actual
event population. Reading the plots: pull narrower than the overlay
$\Rightarrow$ the kinematic-balance constraints add information beyond
the prior alone (Bayesian shrinkage, $\sigma_\text{post}<\sigma_\text{prior}$);
pull matching the overlay shape $\Rightarrow$ prior dominates per event;
pull \emph{wider} than the overlay $\Rightarrow$ unhandled tail or
bias outliers. Title diagnostics: $\langle z\rangle$ (per-event bias)
and $\sigma_z$ (population spread).

\clearpage
\subsection{Detector nuisances}
\begin{figure}[ht]\centering\threecmkin{pulls/kinfit_s1.png}\caption{Jet 1 momentum scale pull.}\end{figure}
\begin{figure}[ht]\centering\threecmkin{pulls/kinfit_s2.png}\caption{Jet 2 momentum scale pull.}\end{figure}
\begin{figure}[ht]\centering\threecmkin{pulls/kinfit_sl.png}\caption{Lepton momentum scale pull.}\end{figure}
\begin{figure}[ht]\centering\threecmkin{pulls/kinfit_sn.png}\caption{MET momentum scale pull.}\end{figure}
\begin{figure}[ht]\centering\threecmkin{pulls/kinfit_t1.png}\caption{Jet 1 $\theta$-shift pull.}\end{figure}
\begin{figure}[ht]\centering\threecmkin{pulls/kinfit_t2.png}\caption{Jet 2 $\theta$-shift pull.}\end{figure}
\begin{figure}[ht]\centering\threecmkin{pulls/kinfit_tl.png}\caption{Lepton $\theta$-shift pull.}\end{figure}
\begin{figure}[ht]\centering\threecmkin{pulls/kinfit_tn.png}\caption{MET $\theta$-shift pull.}\end{figure}
\begin{figure}[ht]\centering\threecmkin{pulls/kinfit_p1.png}\caption{Jet 1 $\phi$-shift pull.}\end{figure}
\begin{figure}[ht]\centering\threecmkin{pulls/kinfit_p2.png}\caption{Jet 2 $\phi$-shift pull.}\end{figure}
\begin{figure}[ht]\centering\threecmkin{pulls/kinfit_pl.png}\caption{Lepton $\phi$-shift pull.}\end{figure}
\begin{figure}[ht]\centering\threecmkin{pulls/kinfit_pn.png}\caption{MET $\phi$-shift pull.}\end{figure}

\clearpage
\subsection{BES, ISR and mass-loss}
\begin{figure}[ht]\centering\threecmkin{pulls/kinfit_bes_m_minus_ecm.png}\caption{BES $\Delta_m$ pull.}\end{figure}
\begin{figure}[ht]\centering\threecmkin{pulls/kinfit_bes_pz.png}\caption{BES $P_z$ pull.}\end{figure}
\begin{figure}[ht]\centering\threecmkin{pulls/kinfit_m_loss.png}\caption{$m(\WW)-m(ee)$ residual pull.}\end{figure}
\begin{figure}[ht]\centering\threecmkin{pulls/kinfit_isr_px.png}\caption{ISR $p_x$ pull.}\end{figure}
\begin{figure}[ht]\centering\threecmkin{pulls/kinfit_isr_py.png}\caption{ISR $p_y$ pull.}\end{figure}
\begin{figure}[ht]\centering\threecmkin{pulls/kinfit_isr_pz.png}\caption{ISR $p_z$ pull.}\end{figure}


\section{Post-fit kinematics: fit vs.\ reco vs.\ gen}\label{app:postfit-kinem}

For each post-fit physical quantity below, the green histogram is the
kinfit value, the orange line is the gen-level truth, and the
light-blue shaded bars are the pre-fit reco values.

\clearpage
\subsection{$W$ mass}
\begin{figure}[ht]\centering\threecmkin{fit_vs_reco_gen/kinfit_Wlep_m.png}\caption{Leptonic $W$ mass.}\end{figure}
\begin{figure}[ht]\centering\threecmkin{fit_vs_reco_gen/kinfit_Whad_m.png}\caption{Hadronic $W$ mass.}\end{figure}

\clearpage
\subsection{Constituent momenta}
\begin{figure}[ht]\centering\threecmkin{fit_vs_reco_gen/kinfit_jet1_p.png}\caption{Jet 1 momentum.}\end{figure}
\begin{figure}[ht]\centering\threecmkin{fit_vs_reco_gen/kinfit_jet2_p.png}\caption{Jet 2 momentum.}\end{figure}
\begin{figure}[ht]\centering\threecmkin{fit_vs_reco_gen/kinfit_lep_p.png}\caption{Lepton momentum.}\end{figure}
\begin{figure}[ht]\centering\threecmkin{fit_vs_reco_gen/kinfit_nu_p.png}\caption{Neutrino momentum (post-fit only at gen level).}\end{figure}

\clearpage
\subsection{\WW system}
\begin{figure}[ht]\centering\threecmkin{fit_vs_reco_gen/kinfit_WW_m.png}\caption{\WW invariant mass.}\end{figure}
\begin{figure}[ht]\centering\threecmkin{fit_vs_reco_gen/kinfit_WW_pz.png}\caption{\WW longitudinal momentum.}\end{figure}


\section{ECM-comparison plots}\label{app:ecm-comparison}

The ECM-comparison plots overlay the three energy points for the same
branch, with the corresponding prior PDF curve in the matching colour.

\begin{figure}[ht]\centering
  \includegraphics[width=0.48\linewidth]{kinfit_vars/ecm_comparison/posteriors/kinfit_bes_m_minus_ecm.png}
  \includegraphics[width=0.48\linewidth]{kinfit_vars/ecm_comparison/posteriors/kinfit_bes_pz.png}
  \caption{Post-fit BES nuisances overlaid for the three energies.}
\end{figure}

\begin{figure}[ht]\centering
  \includegraphics[width=0.48\linewidth]{kinfit_vars/ecm_comparison/fit_vs_reco_gen/kinfit_Wlep_m.png}
  \includegraphics[width=0.48\linewidth]{kinfit_vars/ecm_comparison/fit_vs_reco_gen/kinfit_Whad_m.png}
  \caption{Post-fit leptonic and hadronic $W$ masses overlaid for the
    three energies.}
\end{figure}

\begin{figure}[ht]\centering
  \includegraphics[width=0.48\linewidth]{kinfit_vars/ecm_comparison/fit_vs_reco_gen/kinfit_WW_m.png}
  \includegraphics[width=0.48\linewidth]{kinfit_vars/ecm_comparison/fit_vs_reco_gen/kinfit_WW_pz.png}
  \caption{Post-fit \WW mass and longitudinal momentum overlaid.}
\end{figure}

\end{document}
